% GENERATED by scripts/build-lecture-slides.py; do not edit. Portable profile: XeTeX, theme and fonts from the core.
\makeatletter\IfFileExists{core/beamer/beamerthememx.sty}{\def\mxcorepath{core/}\def\input@path{{core/beamer/}{core/}}}{\def\mxcorepath{../core/}\def\input@path{{../core/beamer/}{../core/}}}\makeatother
\documentclass[aspectratio=169,10pt,t]{beamer}
\usetheme{mx}
\usepackage{graphicx,booktabs,mathtools,array,adjustbox}
\usepackage[all]{xy}
\providecommand{\actualtextwidth}{\textwidth}
\providecommand{\veccor}[2]{\left(\begin{matrix}#1\\#2\end{matrix}\right)}
\DeclarePairedDelimiter{\VecNorm}{\lVert}{\rVert}
\providecommand{\ETC}{\ldots}
\usepackage[french]{babel}
\setlength{\parskip}{.45em}
\hypersetup{pdfauthor={},pdfsubject={Lecture slides},pdfkeywords={},pdftitle={Limites, continuité et dérivation}}
\providecommand{\N}{\mathbb{N}}
\providecommand{\Z}{\mathbb{Z}}
\providecommand{\D}{\mathbb{D}}
\providecommand{\Q}{\mathbb{Q}}
\providecommand{\R}{\mathbb{R}}
\providecommand{\C}{\mathbb{C}}
\providecommand{\vect}[1]{\overrightarrow{#1}}
\providecommand{\vi}{\vec{\imath}}
\providecommand{\vj}{\vec{\jmath}}
\providecommand{\vk}{\vec{k}}
\providecommand{\norm}[1]{\left\lVert #1\right\rVert}
\providecommand{\vc}[2]{\begin{pmatrix}#1\\#2\end{pmatrix}}
\providecommand{\deter}[2]{\det\left(#1,#2\right)}
\providecommand{\abs}[1]{\left\lvert #1\right\rvert}
\providecommand{\dd}{\,\mathrm{d}}
\providecommand{\e}{\mathrm{e}}
\providecommand{\Card}{\operatorname{Card}}
\providecommand{\E}{\operatorname{E}}
\providecommand{\V}{\operatorname{V}}
\providecommand{\degre}{^{\circ}}
\providecommand{\un}[1]{\,\mathrm{#1}}
\providecommand{\sci}[2]{#1\times10^{#2}}
\providecommand{\pgcd}{\operatorname{PGCD}}
\providecommand{\ppcm}{\operatorname{PPCM}}
\title{Limites, continuité et dérivation}
\subtitle{Cours et exemples}
\date{}
\author{}
\begin{document}
\mxtitleframe
\begin{frame}[allowframebreaks]{Limites et continuité}
Les limites aux infinis se calculent de manière similaire à celles des suites.
            
            La limite en \(\pm\infty\) d'une fraction rationnelle est égale à la limite du quotient des termes de plus haut degré.
            example
                  \[
                  \lim_{x\rightarrow+\infty} \dfrac{x^3-23x^2+37}{-2x^2-57} = \lim_{x\rightarrow+\infty} \dfrac{x^3}{-2x^2} 
                                                                           = \lim_{x\rightarrow+\infty} \dfrac{x}{-2}      
                                                                           = -\infty	.
                  \]

example
\end{frame}
\begin{frame}[allowframebreaks]{Limites et continuité}
Soient \(f\) une fonction et \(x_0\) un nombre réel.
\end{frame}
\begin{frame}[allowframebreaks]{Limites et continuité (2)}
\begin{itemize}
\item On dit que \(f\) tend vers \(+\infty\) quand \(x\) tend vers \(x_0\) si 
                        \(f(x)\) devient positif et aussi grand que l'on veut
                        dès que \(x\) est suffisamment proche de \(x_0\).

\end{itemize}
\end{frame}
\begin{frame}[allowframebreaks]{Limites et continuité (3)}
\begin{itemize}
\item On dit que \(f\) tend vers \(-\infty\) quand \(x\) tend vers \(x_0\) si 
                        \(f(x)\) devient négatif et aussi grand que l'on veut en valeur absolue
                        dès que \(x\) est suffisamment proche de \(x_0\).

\end{itemize}
\end{frame}
\begin{frame}[allowframebreaks]{Limites et continuité (4)}
\begin{itemize}
\item On appelle limite de \(f\) à gauche de \(x_0\) 
                     
                        la limite de \(f\) quand \(x\) tend vers \({x_0}\) par valeur inférieure.
                        On note 
                         
                              \(\displaystyle{ \lim_{\substack{x\rightarrow x_0\\x<x_0}} f(x) }\)
                              ou encore \(\displaystyle{ \lim_{x\rightarrow {x_0}^-} f(x) }\).

\end{itemize}
\end{frame}
\begin{frame}[allowframebreaks]{Limites et continuité (5)}
\begin{itemize}
\item On appelle limite de \(f\) à droite de \({x_0}\) 
                     
                        la limite de \(f\) quand \(x\) tend vers \({x_0}\) par valeur supérieure.
                        On note 
                         
                              \(\displaystyle{ \lim_{\substack{x\rightarrow {x_0}\\x>{x_0}}} f(x) }\)
                              ou encore \(\displaystyle{  \lim_{x\rightarrow {x_0}^+} f(x) }\).

\end{itemize}
\end{frame}
\begin{frame}[allowframebreaks]{Limites et continuité (6)}
keypoint 

\begin{itemize}

            \item Il est possible que \(\lim_{x\rightarrow a^+} f(x)\) soit différente de \(\lim_{x\rightarrow a^-} f(x)\).
            
\item Si \(\lim_{x\rightarrow a^+} f(x) = \ell = \lim_{x\rightarrow a^-} f(x)\) 
                  alors \(\lim_{x\rightarrow a} f(x)\) existe et est égale à \(\ell\).

\end{itemize}

      keypoint
\end{frame}
\begin{frame}[allowframebreaks]{Limites et continuité}
\footnotesize
Soit \(f\) une fonction.
            
            Si \(\lim_{x\rightarrow+\infty}f(x)=b\) (ou \(\lim_{x\rightarrow-\infty}f(x)=b\)), avec \(b\) réel,
            la droite d'équation \(y=b\) est alors appelée une \emph{asymptote horizontale}
            à la courbe représentative de \(f\).

            Si \(\lim_{x\rightarrow {a}^{+}}f(x)\) ou \(\lim_{x\rightarrow {a}^{-}}f(x)\) est infinie (\(+\infty\) ou \(-\infty\)),
            la droite d'équation \(x=a\) est alors appelée une \emph{asymptote verticale}
            à la courbe représentative de \(f\).
            
            example 

            example
\end{frame}
\begin{frame}[allowframebreaks]{Limites et continuité}
Les opérations sur les limites se font comme d'habitude :
            table[H]
\end{frame}
\begin{frame}[allowframebreaks]{Limites et continuité (2)}
\footnotesize
\begin{center}\begin{tabular}{lll}
\(\lim f\) & \(\lim g\) & \(\lim (f+g)\) \\ {}
\(\ell\) & \(\ell'\) & \(\ell+\ell'\) \\ {}
\(\ell\) & \(+\infty\) & \(+\infty\) \\ {}
\(\ell\) & \(-\infty\) & \(-\infty\) \\ {}
\(+\infty\) & \(+\infty\) & \(+\infty\) \\ {}
\(-\infty\) & \(-\infty\) & \(-\infty\) \\ {}
\(+\infty\) & \(-\infty\) & Forme indéterminée \\ {}
\end{tabular}\end{center}
\end{frame}
\begin{frame}[allowframebreaks]{Limites et continuité (3)}
Limite d'une somme.
            
            table

table[H]
\end{frame}
\begin{frame}[allowframebreaks]{Limites et continuité (4)}
\footnotesize
\begin{center}\begin{tabular}{lll}
\(\lim f\) & \(\lim g\) & \(\lim (f\cdot g)\) \\ {}
\( \ell \) & \( \ell' \) & \( \ell \cdot \ell' \) \\ {}
\( \ell (\neq 0) \) & \( \infty \) & \(\infty\) (signe du produit) \\ {}
\( \infty \) & \(\infty \) & \(\infty \) (signe du produit) \\ {}
\(0\) & \(\infty\) & Forme indéterminée \\ {}
\end{tabular}\end{center}
\end{frame}
\begin{frame}[allowframebreaks]{Limites et continuité (5)}
Limite d'un produit.
            
            table

table[H]
\end{frame}
\begin{frame}[allowframebreaks]{Limites et continuité (6)}
\footnotesize
\begin{center}\begin{tabular}{lll}
\(\lim f\) & \(\lim g\) & \(\lim \left( \dfrac{f}{g} \right)\) \\ {}
\( \ell \) & \( \ell' (\neq 0) \) & \( \dfrac{\ell}{\ell'} \) \\ {}
\( \ell \) & \(\infty \) & \(0\) \\ {}
\( \ell (\neq 0) \) & \( 0 \) (avec \(g\) de signe constant) & \(\infty\) (règle des signes) \\ {}
\( \infty \) & \( \ell' \) & \(\infty \) (règle des signes) \\ {}
\(0\) & \(0\) & Forme indéterminée \\ {}
\(\infty\) & \(\infty\) & Forme indéterminée \\ {}
\end{tabular}\end{center}
\end{frame}
\begin{frame}[allowframebreaks]{Limites et continuité (7)}
Limite d'un quotient.
            
            table
\end{frame}
\begin{frame}[allowframebreaks]{Limites et continuité}
\footnotesize
Soient \(f\) et \(g\) deux fonctions et soit \(h = g\circ f\), c'est-à-dire \(h(x) = g\left( f(x) \right)\).
            Si \( \lim_{x\rightarrow x_0} f(x) = y_0 \) 
            et \( \lim_{x\rightarrow y_0} g(x) = z_0\) 
            alors \(\displaystyle{\lim_{x\rightarrow x_0} h(x) = z_0} \).
            
            Ici, \(x_0\), \(y_0\) et \(z_0\) peuvent être des nombres réels ou infinis.
            example
                  Soit \(h(x) = \cos \left( \dfrac{1}{x^2 + 1} \right)\).
                  On veut calculer la limite de \(h(x)\) quand \(x\) tend vers \(-\infty\).
                  
                  On pose \(f(x) = \dfrac{1}{x^2 + 1}\) et \(g(x) = \cos x\).
                  Il vient que \(g\left(f(x)\right) = \cos \left( \dfrac{1}{x^2 + 1} \right)\).
                  
                  \[
                  \left.
                  \begin{matrix}
                        \displaystyle{\lim_{x\rightarrow-\infty} f(x) = 0} \\
                        \displaystyle{\lim_{x\rightarrow 0} g(x) = 1} \\
                  \end{matrix}
                  \right\rbrace 
                  \Rightarrow
                  \lim_{x\rightarrow-\infty} h(x) = 1
                  \]
\end{frame}
\begin{frame}[allowframebreaks]{Limites et continuité (2)}
example
\end{frame}
\begin{frame}[allowframebreaks]{Limites et continuité}
\footnotesize
Soit \(f\) une fonction définie sur un intervalle \(I\) et soit \(a \in I\).
            On dit que \(f\) est continue en \(a\) si 
                  \(\lim_{x\rightarrow a^-} f(x) = \lim_{x\rightarrow a^+} f(x) = f(a)\).
            On dit que \(f\) est continue sur \(I\) si elle est continue en tout point de \(I\).
            
            Soit \(f\) une fonction définie à gauche et à droite de \(a\), mais pas en \(a\).
            Si
                  \(\lim_{x\rightarrow a^-} f(x) = \lim_{x\rightarrow a^+} f(x) = \ell\) avec \(\ell\) \emph{réel}, 
            alors on peut prolonger \(f\) par continuité en \(a\) en posant \(f(a) = \lim_{x\rightarrow a} f(x)\).
\end{frame}
\begin{frame}[allowframebreaks]{Limites et continuité}
\begin{itemize}
\item Les fonctions polynômiales ainsi que les fonction \(\sin\) et \(\cos\) sont continues.
            
\item La fonction racine carrée est continue sur \([0,+\infty[\) ; plus généralement, les fonctions usuelles sont continues sur leur ensemble de définition.
            
\item Le produit et la somme de deux fonctions continues sont continus.

\end{itemize}
\end{frame}
\begin{frame}[allowframebreaks]{Limites et continuité (2)}
\begin{itemize}
\item Le quotient de deux fonctions continues est continu sur tout intervalle dans lequel le dénominateur ne s'annule pas.
            
\item La composée de deux fonctions continues est continue.

\end{itemize}
\end{frame}
\begin{frame}[allowframebreaks]{Limites et continuité}
Soit \(f\) une fonction continue sur un intervalle \([a,b]\).
            On note respectivement par \(m\) et \(M\) le minimum et le maximum de \(f\) sur \([a,b]\).
            Alors, pour tout \(y_0\in[m,M]\), il existe \(x_0\in[a,b]\) tel que \(f(x_0)=y_0\).      

En particulier, pour tout réel \(k\) compris entre \(f(a)\) et \(f(b)\), l'équation \(f(x) = k\) admet au moins une solution dans \([a,b]\) ;
            si \(f(a)\) et \(f(b)\) sont de signes contraires, l'équation \(f(x) = 0\) admet au moins une solution dans \([a,b]\).
\end{frame}
\begin{frame}[allowframebreaks]{Limites et continuité}
Soient \(f\), \(g\), \(h\) des fonctions définies au voisinage de \(x_0\) (\(x_0\) réel ou infini).
\end{frame}
\begin{frame}[allowframebreaks]{Limites et continuité (2)}
\begin{itemize}

                  \item Si \(f(x) \leq g(x)\) au voisinage de \(x_0\) et \(\lim_{x_0} f = +\infty\), alors \(\lim_{x_0} g = +\infty\)
                        (de même, si \(g \leq f\) et \(\lim_{x_0} f = -\infty\), alors \(\lim_{x_0} g = -\infty\)).
                  
\item (Théorème des gendarmes) Si \(f(x) \leq g(x) \leq h(x)\) au voisinage de \(x_0\) et
                        \(\lim_{x_0} f = \lim_{x_0} h = \ell\) (réel), alors \(\lim_{x_0} g = \ell\).

\end{itemize}
\end{frame}
\begin{frame}[allowframebreaks]{Limites et continuité (3)}
Par exemple, \(-\dfrac{1}{x} \leq \dfrac{\sin x}{x} \leq \dfrac{1}{x}\) pour \(x > 0\), donc \(\lim_{x\rightarrow+\infty} \dfrac{\sin x}{x} = 0\).
            On admet aussi \(\lim_{x\rightarrow 0} \dfrac{\sin x}{x} = 1\).
\end{frame}
\begin{frame}[allowframebreaks]{Limites et continuité}
Soit \(f : E \rightarrow F\) une fonction.
            On dit que \(f\) est \emph{bijective} si :
\end{frame}
\begin{frame}[allowframebreaks]{Limites et continuité (2)}
\begin{itemize}

                        \item \textbf{Injectivité :}  

                                               si \(a_1\) et \(a_2\) sont des éléments de \(E\) tels que 
                                               \(f(a_1) = f(a_2)\), alors \(a_1=a_2\).
                        
\item \textbf{Surjectivité :} 

                                               pour tout \(b\in F\), il existe \(a\in E\) tel que \(f(a)=b\).

\end{itemize}
\end{frame}
\begin{frame}[allowframebreaks]{Limites et continuité}
Soit \(f\) une fonction continue et strictement monotone sur un intervalle \(I\). 
            Alors \(f\) réalise une bijection de \(I\) vers \(f(I)\).
            
            La fonction \(f\) admet alors une fonction réciproque \(f^{-1} : f(I) \rightarrow I\).
\end{frame}
\begin{frame}[allowframebreaks]{Limites et continuité (2)}
Cette réciproque \(f^{-1}\) est aussi une bijection continue et strictement monotone.
            De plus, \(f\) et \(f^{-1}\) varient dans le même sens.
            Si \(f\) est dérivable en \(x_0 \in I\) avec \(f'(x_0) \neq 0\), alors \(f^{-1}\) est dérivable en \(y_0 = f(x_0)\) et
            \(\left(f^{-1}\right)'(y_0) = \dfrac{1}{f'(x_0)}\) ; si \(f'(x_0) = 0\), la courbe de \(f^{-1}\) a une tangente verticale au point d'abscisse \(y_0\).
\end{frame}
\begin{frame}[allowframebreaks]{Limites et continuité}
Soit \(f\) une fonction continue et strictement monotone sur un intervalle \([a,b]\).
            On note respectivement par \(m\) et \(M\) le minimum et le maximum de \(f\) sur \([a,b]\).
            Alors, pour tout \(y_0\in[m,M]\), il existe \emph{un unique} \(x_0\in[a,b]\) tel que \(f(x_0)=y_0\).
\end{frame}
\begin{frame}[allowframebreaks]{Dérivation}
\footnotesize
Soient \(f\) une fonction et \(x_0 \in \mathbb{R}\).
            
            On dit que la fonction \(f\) est dérivable au point \(x_0\) si 
                     \( \lim_{x\rightarrow x_0} \dfrac{f(x) - f(x_0)}{x-x_0} \)   
            existe et est finie.
            Cette limite est appelée \emph{nombre dérivé} de \(f\) en \(x_0\) et est notée \(f'(x_0)\).
            
            La fonction \(f\) est dite dérivable sur un intervalle \(]a,b[\) si elle est dérivable en tout point de \(]a,b[\).
            L'application \(f' : ]a,b[ \rightarrow \mathbb{R} , \ x \mapsto f'(x)\) est appelée 
            fonction dérivée de \(f\).
            
            Toute fonction dérivable est aussi continue. Mais la réciproque est fausse.
\end{frame}
\begin{frame}[allowframebreaks]{Dérivation}
\footnotesize
Soient \(f\) une fonction et \(x_0 \in \mathbb{R}\).
            On dit que la fonction \(f\) est dérivable à gauche (resp. à droite) en \(x_0\) si la fonction 
                  \(x \mapsto \dfrac{f(x)-f(x_0)}{x-x_0} \)
            admet une limite finie à gauche (resp. à droite).
            Cette limite est alors appelée nombre dérivé de \(f\) à gauche (resp. à droite) en \(x_0\)
            et notée \(f'_g (x_0)\) (resp. \(f'_d (x_0)\) ).
            
            Si \(f\) est une fonction de la variable \(x\), sa dérivée par rapport à \(x\) est aussi notée 
                  \(\dfrac{\mathrm{d}f}{\mathrm{d}x}\).
            
            \(f\) est dérivable en \(x_0\) si et seulement si \(f\) est dérivable à gauche et à droite en \(x_0\) avec \(f'_g(x_0) = f'_d(x_0)\).
            Si \(f'_g(x_0) \neq f'_d(x_0)\), le point de la courbe d'abscisse \(x_0\) est un \emph{point anguleux} ;
            si le taux de variation tend vers l'infini, la courbe a une demi-tangente verticale.
\end{frame}
\begin{frame}[allowframebreaks]{Dérivation}
\footnotesize
Soit \(f\) une fonction dérivable sur un intervalle \(I\).
            Si \(f'\) est aussi dérivable sur \(I\), alors sa dérivée, 
            notée \(f''\) ou \(\dfrac{\mathrm{d}^2f}{\mathrm{d}x^2}\) est appelée
            \emph{dérivée seconde} de \(f\) ; on dit que \(f\) est deux fois dérivable sur \(I\).
            Plus généralement, si \(f\) est \(n\) fois dérivable sur \(I\), 
            on note \(f^{(n)}\) ou \(\dfrac{\mathrm{d}^nf}{\mathrm{d}x^n}\)
            sa dérivée \(n^{\text{ième}}\).
            Si \(f\) admet une dérivée d'ordre \(n\) pour tout \(n\in \mathbb{N}^*\), alors \(f\) est dite
            \emph{indéfiniment dérivable} sur \(I\).
\end{frame}
\begin{frame}[allowframebreaks]{Dérivation}
Soient \(f\) une fonction, 
            \(\mathcal{C}\) la courbe représentative de \(f\) 
            et \(x_0\in \mathbb{R}\)
            tels que \(f'(x_0)\) existe.

Le nombre \(f'(x_0)\) représente alors 
            le coefficient directeur de la droite \((\mathcal{T}_{x_0})\), 
            \emph{tangente} à \(\mathcal{C}\) au point d'abscisse \(x_0\).

L'équation de la droite \((\mathcal{T}_{x_0})\) peut s'écrire
            \( y = f'(x_0) \cdot (x-x_0) + f(x_0) \).
\end{frame}
\begin{frame}[allowframebreaks]{Dérivation}
Si \(f\) est dérivable sur un intervalle \(I\) et si \(f'(x) > 0\) (resp. \(f'(x)<0\)) pour tout \(x \in I\),
                  sauf peut-être en un nombre fini de points où \(f'\) s'annule,
                  alors \(f\) est strictement croissante (resp. strictement décroissante) sur \(I\).
                  Si \(f'(x) = 0\) pour tout \(x \in I\), alors \(f\) est constante sur \(I\).
\end{frame}
\begin{frame}[allowframebreaks]{Dérivation}
Soient \(a\) et \(b\) deux nombres réels tels que \(a<b\).
                        Soit \(f\) une fonction continue sur \([a,b]\) et dérivable sur \(]a,b[\).
                        Alors il existe un réel \(c\in]a,b[\) tel que 
                              
                                    \[
                                          \dfrac{f(b)-f(a)}{b-a}=f'(c)
                                    \]

.
\end{frame}
\begin{frame}[allowframebreaks]{Dérivation}
Soit \(f\) une fonction continue sur \([a,b]\) et dérivable sur \(]a,b[\).
                        Si \(f(a) = f(b)\) alors il existe \(c\in]a,b[\) tel que \(f'(c) = 0\).
\end{frame}
\begin{frame}[allowframebreaks]{Dérivation}
Soit \(f:[a,b]\rightarrow\mathbb{R}\) une fonction (avec \(a\) et \(b\) deux réels tels que \(a<b\)).
                        Si 

\begin{itemize}

                              \item \(f\) est continue sur \([a,b]\) et dérivable sur \(]a,b[\)
                              
\item il existe \(M\in\mathbb{R}\) tel que \(|f'(x)|\leq M\) pour tout \(x\in]a,b[\)

\end{itemize}
\end{frame}
\begin{frame}[allowframebreaks]{Dérivation (2)}
alors \[ \left| \dfrac{f(b) - f(a)}{b - a} \right| \leq M . \]

Plus généralement, si \(|f'(x)| \leq M\) sur un intervalle \(I\), alors \(|f(x) - f(y)| \leq M |x - y|\) pour tous \(x\), \(y\) de \(I\) ;
                        et si \(m \leq f'(x) \leq M\) sur \(]a,b[\), alors \(m(b-a) \leq f(b) - f(a) \leq M(b-a)\).
\end{frame}
\begin{frame}[allowframebreaks]{Dérivation}
Application aux suites \(u_{n+1} = f(u_n)\)

 Soit \(f\) une fonction définie sur un intervalle \(I\) telle que \(f(I) \subset I\), et \((u_n)\) définie par \(u_0 \in I\) et
                        \(u_{n+1} = f(u_n)\) ; alors \(u_n \in I\) pour tout \(n\) (récurrence).
\end{frame}
\begin{frame}[allowframebreaks]{Dérivation (2)}
\begin{itemize}
\item Si \((u_n)\) converge vers \(\ell \in I\) et si \(f\) est continue en \(\ell\), alors \(f(\ell) = \ell\) :
                                    la limite est une solution de l'équation \(f(x) = x\).

\end{itemize}
\end{frame}
\begin{frame}[allowframebreaks]{Dérivation (3)}
\begin{itemize}
\item Si \(f(\ell) = \ell\) avec \(\ell \in I\) et \(|f'(x)| \leq k < 1\) sur \(I\), l'inégalité des accroissements finis donne
                                    \(|u_{n+1} - \ell| \leq k\,|u_n - \ell|\), puis \(|u_n - \ell| \leq k^n |u_0 - \ell|\) ; donc \((u_n)\) converge vers \(\ell\).

\end{itemize}
\end{frame}
\begin{frame}[allowframebreaks]{Dérivation}
table[H]
\end{frame}
\begin{frame}[allowframebreaks]{Dérivation (2)}
\footnotesize
\begin{center}\begin{tabular}{lll}
Fonctions & Dérivées & Domaines de dérivabilités \\ {}
\(c\) & \(0\) & \(\mathbb{R}\) \\ {}
\(\lambda x\) & \(\lambda\) & \(\mathbb{R}\) \\ {}
\(x^n\) (\(n\in\mathbb{N}^*\)) & \(n \cdot x^{n-1}\) & \(\mathbb{R}\) \\ {}
\(\sqrt{x}\) & \(\dfrac{1}{2\sqrt{x}}\) & \(]0 ; +\infty [\) \\ {}
\(\dfrac{1}{x}\) & \(-\dfrac{1}{x^2}\) & \(]-\infty ; 0[ \cup ] 0 ; +\infty[\) \\ {}
\(\sin {x}\) & \(\cos {x}\) & \(\mathbb{R}\) \\ {}
\(\cos {x}\) & \(-\sin {x}\) & \(\mathbb{R}\) \\ {}
\end{tabular}\end{center}
\end{frame}
\begin{frame}[allowframebreaks]{Dérivation (3)}
Dérivée de fonctions usuelles.
      table
\end{frame}
\begin{frame}[allowframebreaks]{Dérivation}
\footnotesize
table[H]
            \begin{center}\begin{tabular}{>{\raggedright\arraybackslash}p{0.29\linewidth}>{\raggedright\arraybackslash}p{0.29\linewidth}>{\raggedright\arraybackslash}p{0.29\linewidth}}
Fonctions & Dérivées & Domaines de dérivabilités \\ {}
\(\lambda u\) & \(\lambda u'\) & en tout réel où \(u\) est dérivable \\ {}
\(u^n\) (\(n\in\mathbb{N}^*\)) & \(n \cdot u' \cdot u^{n-1}\) & en tout réel où \(u\) est dérivable \\ {}
\(\sqrt{u}\) & \(\dfrac{u'}{2\sqrt{u}}\) & tout réel où \(u\) est dérivable et strictement positif \\ {}
\(\dfrac{1}{u}\) & \(-\dfrac{u'}{u^2}\) & tout réel où \(u\) est dérivable et différent de zéro \\ {}
\(\sin {u}\) & \(u' \cos {u}\) & en tout réel où \(u\) est dérivable \\ {}
\(\cos {u}\) & \(- u' \sin {u}\) & en tout réel où \(u\) est dérivable \\ {}
\end{tabular}\end{center}
         Dérivées faisant intervenir des fonctions usuelles.
      table
\end{frame}
\begin{frame}[allowframebreaks]{Dérivation}
table[H]
\end{frame}
\begin{frame}[allowframebreaks]{Dérivation (2)}
\footnotesize
\begin{center}\begin{tabular}{ll}
Fonctions & Dérivées \\ {}
\(u + v\) & \(u' + v'\) \\ {}
\(u\cdot v\) & \(u' \cdot v + u \cdot v'\) \\ {}
\(\displaystyle{ \dfrac{u}{v} }\) & \(\displaystyle{ \dfrac{u' \cdot v - u \cdot v'}{v^2} }\) \\ {}
\(\displaystyle{ f^{-1} }\) & \(\displaystyle{ \dfrac{1}{f'\circ f^{-1}} }\) (là où \(f'\circ f^{-1} \neq 0\)) \\ {}
\(f \circ g\) & \(g' \cdot \left( f' \circ g \right)\) \\ {}
\end{tabular}\end{center}
\end{frame}
\begin{frame}[allowframebreaks]{Dérivation (3)}
Opérations sur les dérivées.
      table
\end{frame}
\begin{frame}[allowframebreaks]{Propriétés géométriques}
Soit \(f\) une fonction et \(\mathcal{D}_f\) son domaine de définition.

\begin{itemize}

                        \item Une fonction \(f\) est dite paire si pour tout \(x\in\mathcal{D}_f\), on a \(-x\in\mathcal{D}_f\) et \(f(-x) = f(x)\).
                        
\item Une fonction \(f\) est dite impaire si pour tout \(x\in\mathcal{D}_f\), on a \(-x\in\mathcal{D}_f\) et \(f(-x) = -f(x)\).

\end{itemize}
\end{frame}
\begin{frame}[allowframebreaks]{Propriétés géométriques}
\textbf{Symétrie par rapport à une droite verticale}

 La courbe représentative d'une fonction paire admet une symétrie par rapport à l'axe \(Oy\).
                  Plus généralement, 
                  la courbe représentative d'une fonction \(f\) admet une symétrie
                  par rapport à la droite d'équation \(x=a\) si
                  \[   \forall x\in\mathcal{D}_f,\ (2a-x)\in\mathcal{D}_f \text{ et } f(2a-x) = f(x)   \]
\end{frame}
\begin{frame}[allowframebreaks]{Propriétés géométriques}
\textbf{Symétrie par rapport à un point}

 La courbe représentative d'une fonction impaire admet une symétrie par rapport à l'origine.
                  Plus généralement, 
                  la courbe représentative d'une fonction \(f\) admet une symétrie par rapport à un point \(I = (a,b)\) si
                  \[   \forall x\in\mathcal{D}_f,\ (2a-x)\in\mathcal{D}_f \text{ et } f(2a-x) + f(x) = 2 b .  \]
\end{frame}
\begin{frame}[allowframebreaks]{Propriétés géométriques}
Une fonction \(f\) est dite périodique de période \(T\) si 
                  \[ \forall x \in \mathcal{D}_f, x+T\in\mathcal{D}_f \text{ et } f(x+T) = f(x)      \;. \]

Les fonctions périodiques les plus fréquemment rencontrées sont 
            les fonctions trigonométriques \(\sin\) et \(\cos\) (de période \(2\pi\)).
\end{frame}
\begin{frame}[allowframebreaks]{Propriétés géométriques (2)}
Pour étudier une fonction périodique de période \(T\), 
            il suffit de l'étudier sur un intervalle de longueur \(T\).
            Pour tracer la courbe, il suffit de le faire sur un tel intervalle,
            le reste de la courbe s'obtient par translation de cette partie.
\end{frame}
\begin{frame}[allowframebreaks]{Propriétés géométriques}
Soient \(I(a,b)\) un point et \(f\) une fonction définie sur un voisinage de \(a\).
            Notons par \((\mathcal{C})\) la courbe représentative de \(f\).
            Le point \(I\) de \((\mathcal{C})\) est un point d'inflexion si 
            la courbe \((\mathcal{C})\) traverse sa tangente au point \(I\) (elle passe d'un côté à l'autre de cette tangente).
\end{frame}
\begin{frame}[allowframebreaks]{Propriétés géométriques (2)}
Si \(f\) est deux fois dérivable au voisinage du point \(I\)
            et si la dérivée seconde s'annule en \(a\) en changeant de signe
            alors \(I\) est un point d'inflexion.
\end{frame}
\begin{frame}[allowframebreaks]{Propriétés géométriques}
Soit \(f : E \rightarrow F\) une fonction bijective, notons par \(\mathcal{C}\) sa courbe représentative.
            Elle admet alors une fonction réciproque \(f^{-1} : F \rightarrow E\) 
            dont on notera par \(\mathcal{C}'\) la courbe représentative.
            Dans un repère orthonormé, les courbes \(\mathcal{C}\) et \(\mathcal{C}'\) sont symétriques par rapport à la première bissectrice (droite d'équation \(y = x\)).
\end{frame}
\begin{frame}[allowframebreaks]{Branches infinies}
Soient \(f\) une fonction et \(\mathcal{C}\) sa courbe représentative.
      La courbe \(\mathcal{C}\) est alors un ensemble de points de la forme \((x,f(x))\).
      On parle de \emph{branche infinie} de la courbe quand au moins une des coordonnées tend vers l'infini.

On sait déjà que si \(\lim_{x\rightarrow x_0} f(x) = \infty\) (à gauche ou à droite) alors 
      la courbe \(\mathcal{C}\) admet une asymptote verticale d'équation \(x=x_0\).

Pour ce qui est de la limite de \(f(x)\) quand \(x\) tend vers l'infini, il y a plusieurs possibilités.
\end{frame}
\begin{frame}[allowframebreaks]{Branches infinies (2)}
Si \(\lim_{x\rightarrow\infty} f(x) = \infty\) et que \(\lim_{x\rightarrow\infty} \dfrac{f(x)}{x} = 0\) alors la courbe \(\mathcal{C}\) admet une branche parabolique de direction \(x'Ox\).
\end{frame}
\begin{frame}[allowframebreaks]{Branches infinies}
Axe \((x'Ox)\)

\[
		\left\{\begin{array}{l}
			\lim_{x\rightarrow \infty} f(x) = \infty 		\\
			\lim_{x\rightarrow \infty} \dfrac{f(x)}{x} = 0 .
		\end{array}\right.
	\]

La courbe \((\mathcal{C})\) admet une branche parabolique de direction \((x'Ox)\)

D'axe \((y'Oy)\)

\[
		\left\{\begin{array}{l}
			\lim_{x\rightarrow \infty} f(x) = \infty 		\\
			\lim_{x\rightarrow \infty} \dfrac{f(x)}{x} = \infty 
		\end{array}\right.
	\]

La courbe \((\mathcal{C})\) admet une branche parabolique de direction \((y'Oy)\)

D'axe \(y = ax\)
\end{frame}
\begin{frame}[allowframebreaks]{Branches infinies (2)}
\[
		\left\{\begin{array}{l}
			\lim_{x\rightarrow \infty} \dfrac{f(x)}{x} = a \neq 0		\\
			\lim_{x\rightarrow \infty} (f(x) - a \, x) = \infty 
		\end{array}\right.
	\]
\end{frame}
\begin{frame}[allowframebreaks]{Branches infinies}
La courbe \((\mathcal{C})\) admet une branche parabolique de direction \((y = a x)\)
\end{frame}
\end{document}
